\section*{الفصل \ref{ch:b3:bioinformatics} --- المعلوماتية الحيوية وتحليل المتتاليات}

\begin{solution}{exo:b3:bioinformatics:1}
الشاملة: المتتاليتان محاذاتان من طرف إلى طرف، وكل بقية في عمود
--- لبروتينين يُعتقد تماثلهما على طولهما كله، مثل
المتماثلات المتعامدة لإنزيم تدبير منزلي. والموضعية: أفضل زوج من
المتتاليتين الجزئيتين درجةً، والباقي مهمل --- لإيجاد نطاق مشترك (نطاق SH2
في بروتينين إشاريين لا صلة بينهما بغير ذلك)، أو مورّثة في
متتالية جينومية طويلة.
\end{solution}

\begin{solution}{exo:b3:bioinformatics:2}
الحدود $0,-1,-2,-3$ في الاتجاهين. السطر A: $1, 0, -1$. والسطر G: $0, 0, -1$.
والسطر C: $-1, -1, 1$. فالمثلى $F(3,3) = 1$: أي AGC فوق AAC بلا فجوات
(مطابقة، وعدم مطابقة، ومطابقة: $1 - 1 + 1 = 1$).
\end{solution}

\begin{solution}{exo:b3:bioinformatics:3}
$s(a,a) = \lambda^{-1}\log\bigl(q_{aa}/p_{a}^{2}\bigr)$. والتربتوفان
نادر ($p_{W} \approx 0.013$)، فاحتمال تحاذي تربتوفانين
مصادفةً، وهو $p_{W}^{2}$، ضئيل، وزوج التربتوفان المحفوظ دليل
تماثل أقوى بكثير من زوج لوسين محفوظ ($p_{L}
\approx 0.1$)؛ فتكون نسبة لوغاريتم الأرجحية أكبر بمقدار موافق.
\end{solution}

\begin{solution}{exo:b3:bioinformatics:4}
القيمة المتوقعة هي عدد المحاذيات التي درجتها لا تقل عن هذه
والتي تُتوقع مصادفةً في بحث هذا الاستعلام في
قاعدة بيانات بهذا الحجم. و$E = 3$ تعني أن ثلاث درجات كهذه تُتوقع
مصادفةً: فالإصابة ليست دليل تماثل (وقد تكون متماثلًا مع ذلك، لكن
البحث لا يستطيع أن يقول).
\end{solution}

\begin{solution}{exo:b3:bioinformatics:5}
$mn = 400\times 2\times 10^{11} = 8\times 10^{13} = 2^{46.2}$. و$E(45) =
2^{1.2} \approx 2.3$؛ و$E(55) = 2^{-8.8} \approx 2\times 10^{-3}$؛
و$E(65) = 2^{-18.8} \approx 2\times 10^{-6}$. و$E = 10^{-3}$ تحتاج $S' =
46.2 + 10.0 = 56$ بتًّا. والاستعلام من \num{40} بقية له $mn$ أصغر عشر
مرات، أي $2^{42.9}$: فيكفيه $53$ بتًّا --- لكن الاستعلام القصير نادرًا
ما يبلغ حتى ذلك.
\end{solution}

\begin{solution}{exo:b3:bioinformatics:6}
المحتويات المعلوماتية: $2$، و$1$، و$0$، و$2 - H$ حيث $H = -(0.7\log_{2}
0.7 + 3\times 0.1\log_{2} 0.1) = 0.36 + 1.00 = 1.36$، أي $0.64$. فالمجموع
$R = 3.64$ بتّ. والمطابقات المصادفة: $9.2\times 10^{6}$ موضع على الطاقين
$\times 2^{-3.64} \approx 7\times 10^{5}$ --- \omterm{def:b3:bioinformatics:motif}{فالنميطة}
تكاد لا تنفع وحدها.
\end{solution}

\begin{solution}{exo:b3:bioinformatics:7}
إقحام عدة بقايا حدث طفري واحد، فلا ينبغي لكلفته
أن تنمو خطيًا مع طوله: فكلفة فتح زائدًا كلفة استطالة
صغيرة تنمذج ذلك. وجزاء الفتح العالي جدًا يفرض
عدم مطابقات حيث تنبغي فجوة ويخطئ \omterm{def:b3:bioinformatics:alignment}{محاذاة} كل ما بعد
إقحام حقيقي؛ والجزاء المنخفض جدًا يبعثر الفجوات في كل مكان، فيطابق
البقايا مصادفةً ويضخّم التطابق.
\end{solution}

\begin{solution}{exo:b3:bioinformatics:8}
ليس بالضرورة: \omterm{def:b3:bioinformatics:alignment}{فالمحاذاة} تغطي قطعة من \num{130} بقية،
وهي سمة نطاق مشترك لا سمة \omterm{def:b3:genomics:comparative}{متماثل متعامد}
محاذًى على طوله. والاختبار: ابحث ببروتين الذبابة رجوعًا في
مجموع البروتينات البشرية (هل الاستعلام أفضل إصابة له، على الطول كله؟)،
وحدّد النطاق بنموذج ماركوف مخفي ملمحي، وابنِ شجرة مورّثية للأسرة
في عدة أنواع.
\end{solution}

\begin{solution}{exo:b3:bioinformatics:9}
\omterm{def:b3:bioinformatics:msa}{الملمح} يسجّل، عمودًا عمودًا، ما تتحمله الأسرة: موضع
دائم الكراهة للماء وإن لم تكن البقية نفسها، وبقية حفّازة
ثابتة، وموضع دائم الفجوة في نصف
الأسرة. أما \omterm{def:b3:bioinformatics:alignment}{المحاذاة} الثنائية فتدرّج كل بقية في مواجهة بقية
واحدة أخرى ولا تستطيع أن تعلم أن فالينًا عند الموضع 40 ``بجودة''
الإيزولوسين الذي عنده في الاستعلام. كما يزن \omterm{def:b3:bioinformatics:msa}{الملمح}
الأعمدة المحفوظة، فيصير التشابه الضعيف المتركز حيث تُحفظ
الأسرة ذا دلالة.
\end{solution}

\begin{solution}{exo:b3:bioinformatics:10}
عند درجة متوقعة موجبة $\mu > 0$ لكل عمود، تكون الدرجة التراكمية
على امتداد قطر متتاليتين عشوائيتين مسيرًا عشوائيًا بانجراف
موجب: فبعد $n$ عمود تكون نحو $\mu n$، وتكون أفضل
\omterm{def:b3:bioinformatics:alignment}{محاذاة موضعية} هي الكل عمليًا وتنمو درجتها مثل
$\mu n$ لا مثل $\log n$. ونظرية كارلن--ألتشول، التي
تقتضي انجرافًا سالبًا حتى تكون الدرجات العالية جولات نادرة،
لا تنطبق ولا يوجد $\lambda$. أما في DNA عند \qty{60}{\%} من G وC فإن
احتمال المطابقة $2(0.3^{2}) + 2(0.2^{2}) = 0.26$، فتكون \omterm{prop:b3:bioinformatics:logodds}{الدرجة
المتوقعة} $0.26 - 0.74 = -0.48$: أي ما تزال سالبة، والإحصاءات
صالحة؛ لكن نظامًا مثل مطابقة $+1$ وعدم مطابقة $-0.3$ يكون توقعه
$+0.04$ ويبلّغ عن الجينوم كله محاذاةً واحدة.
\end{solution}

\begin{solution}{exo:b3:bioinformatics:11}
البذور والتسلسل: جد مطابقات تامة أو شبه تامة لقطع بطول $k$ بين
المتتاليتين بجدول تجزئة، واحتفظ بما يصطف منها على
أقطار متسقة، وسلسلها، وشغّل \omterm{thm:b3:bioinformatics:nw}{البرمجة الديناميكية} في
الفجوات بين البذور المسلسلة فقط؛ وهو يتخلى عن المحاذيات في مناطق
لا \omterm{def:b3:bioinformatics:blast}{بذرة} فيها (القطع الشديدة التباعد). والتنطيق: إذا كانت المتتاليتان
معلومتي شبه التوازي، فاحسب الخلايا الواقعة داخل نطاق
عرضه $w$ حول القطر فقط، بكلفة $wn$ بدل $mn$؛ وهو يتخلى
عن أي \omterm{def:b3:bioinformatics:alignment}{محاذاة} فيها إقحام أكبر من النطاق.
\end{solution}

\begin{solution}{exo:b3:bioinformatics:12}
التسلسل المشفِّر له دورية ثلاثية: فمواضع الرامزة الثلاثة
تختلف تراكيب قواعدها (والثالث أشدها انحيازًا)، وانحياز
الرامزات غير متساوٍ في كل نوع. والنموذج ذو الحالات المشفِّرة الثلاث
بالتتابع، تصدر كل منها قواعد بتركيب ذلك الموضع من
الرامزة، يسند إلى DNA المشفِّر احتمالًا أعلى مما يسنده
إليه حالة غير المشفِّر، على نافذة من بضع عشرات من الرامزات، حتى بلا وقفات.
وفي الجينوم البشري تكون الإكسونات قصيرة (\qty{150}{bp})، تفصلها إنترونات
بالكيلوقواعد، فتكون الإشارة المشفِّرة قصيرة متقطعة؛ وعلى النموذج
أن يتعرف مواقع التشذيب أيضًا، وهي إشارات ضعيفة، كما أن
كثرة التسلسل غير المشفِّر تنتج قطعًا مشفِّرة كاذبة
كثيرة.
\end{solution}

\begin{solution}{pb:b3:bioinformatics:1}
\textbf{1.} السطور K وQ وT؛ والأعمدة K وA وQ وT؛ والحدود $0,-1,-2,-3,-4$
و$0,-1,-2,-3$. السطر K: $1, 0, -1, -2$؛ والسطر Q: $0, 0, 1, 0$؛ والسطر T:
$-1, -1, 0, 2$. فالمثلى $2$: \texttt{K-QT} فوق \texttt{KAQT}.
\textbf{2.} أفضل درجة موضعية $3$: \texttt{CAT} في مواجهة \texttt{CAT}
(البقايا 4--6 من GATCAT مع 2--4 من ACAT)؛ و\texttt{ATCAT} في مواجهة
\texttt{A-CAT} تعطي أيضًا $4 - 1 = 3$.
\textbf{3.} $300\times 450 = 1.35\times 10^{5}$ تحديث؛ وفي مواجهة
قاعدة البيانات $300\times 1.2\times 10^{11} = 3.6\times 10^{13}$.
\textbf{4.} $3.6\times 10^{4}$ ثانية، أي عشر ساعات لكل استعلام؛ وبذور \omterm{def:b3:bioinformatics:blast}{BLAST}
تتخطى الجدول كله تقريبًا وتجيب في ثوانٍ.
\textbf{5.} $q_{ab} = p_{a}p_{b}e^{\lambda s}$. الألانين: $e^{1.39} = 4.0$،
و$q_{AA} = 0.074^{2}\times 4.0 = 0.022$، والنسبة $4$. والتربتوفان:
$e^{3.82} = 45$، و$q_{WW} = 0.013^{2}\times 45 = 0.0077$، والنسبة $45$. فزوج
التربتوفان المتحاذي أكثر تواترًا في المتماثلات بمقدار $45$ ضعفًا مما هو
مصادفةً، وزوج الألانين بأربعة أضعاف فقط؛ ومع ذلك فأزواج الألانين
أشيع بالقيمة المطلقة لأن الألانين شائع.
\textbf{6.} \qty{24}{\%} تقع في منطقة الشفق، حيث تبلغ المحاذيات
العشوائية \qtyrange{15}{20}{\%}؛ وتحسمها القيمة المتوقعة \omterm{def:b3:bioinformatics:alignment}{للمحاذاة}،
أو نميطات محفوظة في المواضع الصحيحة، أو مطابقة نموذج ماركوف مخفي ملمحي لأسرة
معروفة، أو طية مشتركة.
\textbf{7.} $mn = 300\times 1.2\times 10^{11} = 3.6\times 10^{13}$؛
و$\log_{2}(mn) = 45.0$.
\textbf{8.} $E(92) = 2^{45 - 92} = 2^{-47} \approx 7\times 10^{-15}$؛
و$E(38) = 2^{7} = 128$.
\textbf{9.} $E = 10^{-3}$ عند $S' = 45 + 10 = 55$ بتًّا؛ و$E = 1$ عند $45$
بتًّا.
\textbf{10.} عشرة أضعاف $mn$: أي $E \approx 7\times 10^{-14}$، وهي ما تزال
ساحقة.
\textbf{11.} عند $E = 128$ تكون الإصابة العاشرة ما تنتجه المصادفة؛
وتطابق \qty{40}{\%} على $60$ بقية ليس دليلًا. وبحث بالملامح
للبقايا 200--260 في قاعدة بيانات النطاقات قد يبيّن هل
تلك القطعة نطاق معروف، بإحصاءات تفتقر إليها المقارنة
الثنائية.
\textbf{12.} الإصابات الأفضل المتبادلة على الطول الكامل تتسق
مع تعامد واحد لواحد؛ لكنها لا تثبته --- فمضاعفة في
سلالة واحدة بعد الانفصال تعطي متعامدين شريكين، وفقدُ
المتعامد الحقيقي قد يترك متماثلًا متوازيًا أفضلَ إصابة. والمحك شجرة مورّثية بعدة
أنواع.
\textbf{13.} $R = 2 + 2 + 1.6 + 2 + 0.8 + 1.2 + 0.4 + 0.3 = 10.3$ بتّ.
\textbf{14.} $3.2\times 10^{9}\times 2^{-10.3} \approx 2.5\times 10^{6}$
مطابقة مصادفة.
\textbf{15.} $\log_{2}(3.2\times 10^{9}/200) = \log_{2}(1.6\times 10^{7})
\approx 24$ بتًّا.
\textbf{16.} الورود المشترك عند مباعدة ثابتة يجمع البتّات: $10.3 + 8 =
18.3$، فيسدّ $8$ بتّات من $13.7$ الناقصة؛ ويجب أن يأتي نحو $5.7$ بتّ (أي عامل
$50$ في المطابقات المصادفة) من مكان آخر --- من نفاذية
\omterm{def:b3:chromatin-epigenetics:nucleosome}{الصبغين}، ومن شركاء آخرين.
\textbf{17.} $H = -(0.5\log_{2}0.5 + 0.5\log_{2}0.5) = 1$ بتّ؛ و$R = 2 -
1 = 1$ بتّ.
\textbf{18.} على العامل البكتيري أن يجد مواقعه القليلة في جينوم
\qty{4.6}{Mb} بلا عون من \omterm{def:b3:chromatin-epigenetics:nucleosome}{الصبغين}: فيحتاج نحو $19$
بتًّا، وتكون مواقعه طويلة محفوظة. أما جينوم حقيقيات النوى فأكبر
بألف مرة، فيحتاج عشرة بتّات أخرى، ومع ذلك فمواقع عوامله
قصيرة؛ وهي تبلغ النوعية بالتوليف وبحصر
\omterm{def:b3:chromatin-epigenetics:nucleosome}{الصبغين} النفوذ، وهذا أيضًا يجعل التنظيم أقبل
للتطور، لأن الموقع القصير يسهل كسبه وفقده.
\textbf{19.} $3/64 = 0.047$؛ والعدد المتوقع من الرامزات قبل
وقف هو $64/3 \approx 21$.
\textbf{20.} $p_{A} = p_{T} = 0.31$، و$p_{G} = p_{C} = 0.19$: ومنه $P(\text{TAA})
= 0.31^{3} = 0.030$، و$P(\text{TAG}) = P(\text{TGA}) = 0.31^{2}\times 0.19
= 0.018$؛ والمجموع $0.066$، والطول المتوقع للإطار $15$ رامزة. وDNA الغني
بالقاعدتين A وT مليء بالوقفات، فتكون الأطر المفتوحة العشوائية أقصر والطويلة
منها أبرز.
\textbf{21.} $(61/64)^{100} = e^{-4.80} = 0.008$؛ و$(61/64)^{300} =
e^{-14.4} = 5.6\times 10^{-7}$.
\textbf{22.} كل إطار مفتوح أعظمي ينتهي بوقف، و$3.2\times
10^{9}$ بداية رامزة تحوي $3.2\times 10^{9}\times 3/64 = 1.5\times
10^{8}$ وقفة: أي نحو $1.5\times 10^{8}\times 0.008 = 1.2\times 10^{6}$
إطار عشوائي بطول $100$ رامزة على الأقل، و$1.5\times 10^{8}\times
5.6\times 10^{-7} \approx 80$ بطول $300$ على الأقل.
\textbf{23.} للبكتيريا البالغة \qty{4.6}{Mb} نحو $4\times 10^{5}$ وقفة
ومن ثم نحو \num{3500} إطار مصادفة بطول $100$ رامزة لكن لا يكاد يوجد
شيء بطول $300$؛ ومتوسط مورّثاتها $300$ رامزة و\qty{88}{\%} من DNA فيها
مشفِّر، فالإطار المفتوح الطويل مورّثة دائمًا تقريبًا. أما في الدودة فإن
\qty{1.5}{\%} من DNA مشفِّر، ومتوسط الإكسونات $50$ رامزة --- أقصر
من عتبة المصادفة --- ويغرقها مليون إطار عشوائي بطول $100$
رامزة. وتستعمل كواشف المورّثات في حقيقيات النوى إشارات مواقع التشذيب،
وانحياز الرامزات في نموذج ماركوف مخفي، والتماثل مع بروتينات معروفة، وقبل
ذلك كله النسخ المسلسلة.
\textbf{24.} القراءة الآتية من رسالة مشذَّبة تُحاذى إلى الجينوم في
قطعتين يفصلهما إنترون: فالشق يعلّم موقعي التشذيب
حتى القاعدة؛ وتغطية القراءات ترسم حدود الإكسونات والقراءات المزدوجة تصل
الإكسونات المتعاقبة في نسخة واحدة، فتحلّ أي مواقع تشذيب مرشحة
استُعملت.
\textbf{25.} $E \approx 7\times 10^{-15}$ لأفضل إصابة؛ ونحو $24$
بتًّا لتحديد $200$ موقع في الجينوم؛ ونحو $80$ إطار قراءة مصادفًا
بطول $300$ رامزة في الجينوم كله.
\end{solution}
